\documentclass[11pt,a4paper]{article}
\usepackage[UTF8]{ctex}
\usepackage{geometry}
\geometry{left=2.5cm,right=2.5cm,top=2.5cm,bottom=2.5cm}
\usepackage{amsmath,amssymb,amsthm}
\usepackage{hyperref}
\usepackage{booktabs}
\usepackage{enumitem}
\usepackage{xltabular}
\hypersetup{colorlinks=true,linkcolor=blue,urlcolor=blue}

\newtheorem{definition}{定义}[section]
\newtheorem{proposition}{命题}[section]
\newtheorem{remark}{注记}[section]

\title{AI化攻防：当大模型加入博弈（修正版）\\
——推理成本、对抗样本与柯克霍夫原则的边界条件}
\author{李龙胜}
\date{\today}

\begin{document}

\maketitle

\begin{center}
\textit{
当攻防双方都部署大模型时，成本结构发生质变。\\
推理成本替代人工成本成为主要消耗源，\\
训练成本成为新的战略性投入，\\
对抗样本和对抗训练形成独立的模型参数空间子博弈。\\
贪心择优从“分析还是忽略”扩展到三个新维度：\\
模型选择、推理深度和对抗训练频率。\\
但柯克霍夫原则面临着AI时代的新挑战——\\
当攻击方探测成本趋近于零时，周期性扰动是否仍然足够？\\
本文在全面论述AI化攻防后，专门讨论这一边界条件。
}
\end{center}

\vspace{5mm}

\begin{abstract}
本文建立AI化攻防的理论框架，
讨论当攻防双方都部署大模型时，
成本函数如何变化，贪心择优如何扩展到新维度。
防御方的大模型推理成本、训练与微调的成本摊销、
攻击方的对抗样本生成成本、模型窃取与数据投毒的防御——
全部被纳入统一的贪心择优调度框架。
贪心择优在模型选择、推理深度和对抗训练频率三个新维度上
自然扩展，每个维度的最优解结构与已有框架完全同构。
进一步指出攻击方按需生成与防御方全量分析的成本不对称性，
AI成本陷阱，以及柯克霍夫原则在AI时代的边界条件。
对攻击方大模型能力的对称性假设和无限弹药推论进行边界标注。
\end{abstract}

\tableofcontents

\section{引言}

在已有的安全运营理论框架中，
分析能力主要体现在分析师工时和传统自动化工具的算力消耗上。
随着大模型在攻防双方的广泛应用，
分析能力消耗的形式发生了质变——
防御方使用大模型分析告警和生成检测规则，
攻击方使用大模型生成攻击payload和对抗样本。

本文建立AI化攻防的理论框架，
将大模型的推理成本、训练成本摊销、
对抗样本攻防经济学和模型保护机制
纳入已有的贪心择优调度体系。
在此基础上，专门讨论AI时代对已有框架中多个核心假设的冲击。

\section{大模型推理成本作为分析能力消耗源}

防御方部署大模型分析告警，每次推理消耗Token或GPU时间。
与分析师相比，大模型的成本结构有本质区别：
大模型边际成本低但固定成本高，
单次推理成本远低于人工研判但前期训练投入巨大；
大模型可并行处理多条告警，不受人类注意力瓶颈限制；
但大模型推理质量参差不齐，低价值任务使用最强模型会造成浪费。

最优模型选择由够用就行原则裁定。
设防御方有两种模型可用：
全量模型 \(M_D\)，单次推理成本为 \(c_D^{\text{infer}}\)，检测精度为 \(q_{\text{full}}\)；
轻量级模型 \(M_D^{\text{light}}\)，推理成本仅为 \(c_D^{\text{infer}}/k\)，检测精度为 \(q_{\text{light}} < q_{\text{full}}\)。

对于高A值告警类别，精度损失代价极大——
漏报一条APT告警损失可能达千万级。
使用全量模型推理，\(q_{\text{full}}\) 的精度优势产生的边际收益远超额外推理成本。
对于低A值告警类别，精度损失代价小——
漏报一条扫描告警损失仅百元级。
使用轻量级模型或传统规则引擎即可，
全量模型的额外推理成本大于精度优势带来的收益。
对于中等A值告警类别，仅在UEBA偏离度高时使用全量模型，
低偏离度时使用轻量级模型兜底。

这种分层推理的数学结构与沙箱分析深度的分层优化、
UEBA基线粒度的分层调整完全同构——
高价值任务用高成本高精度工具，低价值任务用低成本低精度工具，
由贪心择优统一裁定最优分配。

\textbf{技术注释：推理成本的预算归属问题}。
大模型推理消耗GPU算力，人工研判消耗分析师工时，
两者在物理上不共享同一资源池，
但在分析能力预算的约束下互相挤占。
在工程实现中，分析能力预算可能需要拆分为多个资源子池——
GPU算力池、人工工时池、存储池——
而贪心择优需要在多个资源约束下做联合调度。
当前版本使用单一分析能力预算作为统一抽象，
在AI化场景中是简化假设。

\section{模型训练与微调的成本摊销}

防御方为特定告警类别微调专用检测模型需要巨大的前期训练投入，
攻击方训练生成式攻击模型同样如此。
训练成本必须分摊到每次推理中，直接影响最优分析率的计算。

设防御方为告警类别 \(i\) 微调了专用检测模型，训练成本为 \(C_i^{\text{train}}\)。
该模型预计在 \(N_i^{\text{expected}}\) 次推理后被下一轮微调替代。
每次推理的有效成本为
\[
c_{\text{eff}} = c_{\text{infer}} + \frac{C_i^{\text{train}}}{N_i^{\text{expected}}}.
\]

若某类告警的日均到达量很低，训练成本的摊销周期极长，
每次推理的有效成本极高——为低频率告警类别微调专用模型不划算，
该类别应使用通用模型或传统规则引擎。
若某类告警的日均到达量极高，训练成本被大量推理分摊，
每次推理的有效成本接近纯推理成本——高频告警类别值得投入专用微调。

最优训练投入由贪心择优裁定：
微调带来的检测精度提升是否足以覆盖微调成本在推理中的摊销。
这恰好是够用就行原则在模型训练维度的直接应用。

\section{对抗样本的攻防经济学}

攻击方制造对抗样本来逃避检测，
防御方进行对抗训练来增强模型鲁棒性。
两者在模型参数空间形成了一个独立的子博弈。

设攻击方生成一个对抗样本的成本为 \(c_A^{\text{adv}}\)，
该样本逃避检测的概率为 \(p_{\text{evade}}\)。
防御方进行一轮对抗训练的成本为 \(c_D^{\text{adv\_train}}\)，
训练后模型对对抗样本的检测率提升 \(\Delta q_{\text{adv}}\)。

攻击方的最优对抗样本投入：
当 \(p_{\text{evade}} \cdot A_{\text{target}} > c_A^{\text{adv}}\) 时，
发起对抗样本攻击。
防御方的最优对抗训练投入：
当 \(\Delta q_{\text{adv}} \cdot A_{\text{target}} \cdot N_{\text{adv\_alerts}} > c_D^{\text{adv\_train}}\) 时，
执行对抗训练。

双方在模型参数空间形成一个新的均衡——
攻击方的对抗样本生成成本和防御方的对抗训练成本互相驱动，
最终在边际收益相等的点上稳定下来。
这与已建立的逼移策略和探测反探测博弈完全同构，
只是作用在模型参数空间而非网络流量和告警流上。

\section{数据投毒与模型窃取的防御}

攻击方可能污染训练数据或通过多次查询来窃取防御方模型。
这部分与已有的柯克霍夫原则自然衔接——
模型权重和训练数据成为新的秘密参数，受柯克霍夫原则保护。

防御方的模型权重和微调数据集属于秘密层保护对象。
模型微调对应周期性扰动——
每个周期用新的对抗样本和数据增强重新微调模型，
使攻击方在上一个周期通过查询推断的模型边界过期。
模型输出上叠加随机噪声，
对抗攻击方通过多次查询来窃取模型边界的探测行为。

这与UEBA基线参数的柯克霍夫保护完全同构——
周期性扰动对抗跨周期探测，随机噪声对抗单周期内试探。

\section{贪心择优的AI维度扩展}

在AI化的攻防中，每一步决策不再是简单的“分析还是忽略”，
而是涉及三个新维度。

\textbf{模型选择维度}：
用全量模型、轻量级模型还是传统规则引擎。
最优选择由告警的A值和当前GPU算力预算的饱和度决定。
当GPU算力充裕时，更多告警使用全量模型；
当GPU算力紧张时，降级为轻量级模型。

\textbf{推理深度维度}：
单轮推理还是多轮思维链推理。
多轮推理精度更高但成本倍增。
最优推理深度由告警复杂度和A值决定——
高A值复杂告警使用多轮推理，低A值简单告警单轮推理即可。

\textbf{对抗训练频率维度}：
每批次、每日还是每周进行一次对抗训练。
频率越高模型鲁棒性越强，但训练算力消耗越大。
最优频率是边际鲁棒性收益等于边际训练成本的均衡点。

三个维度的最优解结构完全一致——
都是边际收益等于边际成本的均衡点。
贪心择优定理在AI维度上同样成立。

\section{攻击方的AI成本优势与防御方的对称性假设}

攻击方用AI是“按需生成”，防御方是“全量分析”，
两者成本结构有本质差异。

设攻击方生成一条能逃避检测的攻击payload的推理成本为 \(c_A^{\text{gen}}\)，
该payload成功绕过防御方AI检测的概率为 \(p_{\text{evade}}\)，
成功后的收益为 \(A_{\text{target}}\)。
攻击方仅在 \(p_{\text{evade}} \cdot A_{\text{target}} > c_A^{\text{gen}}\) 时
使用AI生成模型，否则使用传统攻击方法。

防御方用AI分析一条告警的推理成本 \(c_D^{\text{analyze}}\) 即使很低，
也必须乘以海量告警基数 \(N_{\text{alerts}}\)，
总成本 \(N_{\text{alerts}} \cdot c_D^{\text{analyze}}\) 仍然巨大。

两者的有效成本处于不同的量级——
攻击方按需生成，防御方全量分析。
这意味着在AI化的第一阶段，
攻击方可能比防御方更快地获得收益。

\textbf{对称性假设的边界条件}：
攻击方大模型能力的对称性假设在定性分析中成立——
博弈结构对称，双方都受够用就行原则约束。
但在当前真实攻防态势中，双方的大模型部署成熟度可能存在阶段性不对称。
防御方的大模型部署是可控的、可集成的；
攻击方的大模型部署需要在对抗环境中隐蔽运行，
受限于算力和通信约束。
本文的均衡分析描述的是双方AI能力都成熟后的极限状态。

\section{防御方的“AI成本陷阱”}

防御方引入AI时面临三重陷阱。

\textbf{AI虚警消耗分析能力}。
AI产生误报时，分析师的研判精力被无效告警消耗。
若AI误报率过高，净分析能力收益可能为负——
AI节省的时间被自身误报吞噬殆尽。
这与UEBA误报负反馈循环完全同构。

\textbf{AI模型退化}。
攻击方持续对抗训练导致防御方模型精度缓慢下降。
防御方被迫持续投入再训练成本，
形成“不训练就退化，训练就消耗算力”的两难困境。
最优再训练频率是边际精度维护收益等于边际训练成本的均衡点。

\textbf{AI依赖削弱人的能力}。
分析师长期依赖AI裁定后，自身研判技能退化。
当AI模型因对抗攻击突然失效时，
分析师的研判能力不足以承接。
防御方需要刻意保留一定比例的人工研判，
维持分析师基本研判能力——
即使短期内不最大化AI使用效率。
这是够用就行原则在AI时代的新维度：
不是最大化AI的使用，而是保留必要的人的能力作为后备。

\section{AI化攻防的核心结论与柯克霍夫原则的边界条件}

当攻防双方都部署AI时，优势不会自动归于防御方。
攻击方的AI使用成本更低——生成攻击比分析海量告警更省算力，
攻击方可针对性按需生成，防御方被海量告警基数锁死。
AI化攻防没有消除博弈，只是将战场从网络流量层转移到了模型参数空间，
结构在数学上完全保留。
防御方面临AI成本陷阱——虚警消耗分析能力、模型退化、分析师能力丧失，
三者叠加可能导致总安全能力不升反降。

\textbf{柯克霍夫原则在AI时代的边界条件}：
在第2卷和第3卷中，防御方通过柯克霍夫原则获得不对称优势——
参数保密和周期性扰动让攻击方的探测努力持续过期。
但在AI时代，如果攻击方可以通过大模型无限生成对抗样本、
无限探测防御方模型边界，
攻击方的探测成本可能趋近于零。

周期性扰动对抗无限探测的有效性，
依赖于防御方的扰动频率是否能超过攻击方的探测生成速度。
当攻击方能够以趋近于零的成本持续生成探测时，
周期性扰动可能不足以消耗攻击方的探测预算——
攻击方可以承受所有探测努力的过期。
这不是柯克霍夫原则的失效，
而是它的有效性边界需要在AI时代被重新审视。
这一边界条件的确切形式和极限情况
需要进一步的数值模拟和实战验证。

\section{诚实标注}

\begin{center}
\begin{xltabular}{\textwidth}{c|X|c}
\toprule
\textbf{命题} & \textbf{状态} & \textbf{层级} \\
\midrule
大模型推理成本作为分析能力消耗源 &
成本结构模型已建立，与已有分析能力抽象衔接 &
II（推理成本的具体数值需实际部署标定） \\
分析能力预算的多资源池拆分 &
单一预算在AI化场景中是简化假设，
工程实现中需拆分为GPU算力池、人工工时池等 &
II（技术注释，工程实施时细化） \\
模型训练与微调的成本摊销 &
摊销公式已给出，与够用就行原则对接 &
II（训练成本的具体分摊周期需实际部署标定） \\
对抗样本攻防经济学 &
攻防双方的成本收益模型已建立，与逼移博弈同构 &
II（对抗样本生成和检测的精确成本需实测） \\
攻击方大模型能力的对称性假设 &
定性分析中成立，当前真实攻防态势中双方部署成熟度可能不对称 &
II（推测性分析，描述双方AI能力成熟后的极限状态） \\
数据投毒与模型窃取的防御 &
与柯克霍夫原则自然衔接 &
II \\
贪心择优的AI维度扩展 &
模型选择、推理深度、对抗训练频率三个新维度的最优解结构已明确 &
II \\
防御方的AI成本陷阱——虚警、退化、能力丧失 &
三项陷阱的逻辑已建立 &
II（实战数据待验证） \\
AI化攻防核心结论 &
定性结论已明确，精确均衡解需进一步数值模拟 &
II \\
柯克霍夫原则在AI时代的边界条件 &
周期性扰动对抗无限探测的有效性依赖于扰动频率与探测生成速度之比 &
III（推测性分析，需进一步验证） \\
\bottomrule
\end{xltabular}
\end{center}

\section{结语}

AI化攻防专题卷的初步框架已经建立，并经外部审查修正。
大模型推理成本是新的分析能力消耗源，
模型训练是新的参数重标定维度，
对抗样本攻防是新的逼移子博弈，
模型权重保护是柯克霍夫原则在新参数空间的应用。
贪心择优在模型选择、推理深度和对抗训练频率三个新维度上自然扩展，
每个维度的最优解结构与已有框架完全同构。

柯克霍夫原则在AI时代面临新的边界条件——
攻击方探测成本趋近于零时，
周期性扰动的有效性需要重新验证。
这与第2卷、第3卷的核心推论形成深刻的张力，
值得后续深入研究。

\vspace{5mm}
\noindent\textbf{文档信息}：
\begin{itemize}
    \item 作者：李龙胜
    \item 版本：修正版（整合外部审查反馈）
    \item 许可：CC BY 4.0
    \item 前置依赖：四卷体系、产品蓝图修正版、
          安全参数自适应调节文档、UEBA桥接、
          SOAR剧本优化、扩展分析路径统一调度
\end{itemize}

\end{document}